% Part: first-order-logic
% Chapter: beyond
% Section: other-logics

\documentclass[../../../include/open-logic-section]{subfiles}

\begin{document}

\olfileid[pt]{fol}{byd}{oth}

\olsection{Outras Lógicas}

Como você já deve ter percebido, não é difícil projetar uma nova lógica.
Você também pode criar a sua própria sintaxe, montar um sistema dedutivo e
elaborar uma semântica para acompanhá-la. Talvez você tenha que ser um pouco
esperto se quiser que o sistema de !!{derivation} seja completo para a
semântica, e pode exigir algum esforço convencer o mundo em geral de que a
sua lógica é realmente interessante. Mas, em troca, você pode desfrutar de
horas de diversão sadia, explorando as propriedades matemáticas e
computacionais da sua lógica.

As últimas décadas testemunharam uma verdadeira explosão de lógicas formais.
A lógica difusa é projetada para modelar o raciocínio sobre propriedades
vagas. A lógica probabilística é projetada para modelar o raciocínio sobre a
incerteza. As lógicas \emph{default} (de exceções) e as lógicas não
monotônicas são projetadas para modelar formas revogáveis de raciocínio,
isto é, inferências ``razoáveis'' que podem mais tarde ser derrubadas diante
de novas informações. Há lógicas epistêmicas, projetadas para modelar o
raciocínio sobre o conhecimento; lógicas causais, projetadas para modelar o
raciocínio sobre relações causais; e até mesmo lógicas ``deônticas'', que
são projetadas para modelar o raciocínio sobre obrigações morais e éticas.
Dependendo de a principal motivação para a introdução desses sistemas ser
filosófica, matemática ou computacional, você pode encontrar estudos sobre
essas criaturas sob a rubrica de lógica matemática, lógica filosófica,
inteligência artificial, ciência cognitiva ou outro lugar.

A lista não tem fim, e as possibilidades parecem infinitas. Talvez nunca
alcancemos o sonho de Leibniz de reduzir toda a razão humana a cálculo ---
mas isso não pode nos impedir de tentar.

\end{document}
